\chapter{Justificación}

El desarrollo de experiencias inmersivas en realidad virtual (VR, 
por sus siglas en inglés) ha evolucionado desde esquemas 
centrados únicamente en la interacción básica hacia enfoques que
consideran el comportamiento del usuario como componente para 
modificar la experiencia en tiempo real \cite{lopes2025closing}. 
Estudios recientes muestran que el diseño de interfaces de 
interacción en VR, las técnicas de interacción en entornos 
inmersivos y las interfaces tridimensionales constituyen áreas
activas de investigación, debido a que la calidad de la 
experiencia depende no solo de la representación visual del 
entorno, sino también de la forma en la que el usuario interactúa 
con él \cite{chen2024interfaces,yeo2024interfaces}. Bajo este 
enfoque, es importante plantear la creación de una experiencia 
que no se limite a implementar mecánicas interactivas, sino que 
incorpore una capa adaptativa capaz de responder al 
comportamiento del jugador.

La relevancia de este proyecto se encuentra en la necesidad de 
avanzar hacia experiencias interactivas más sensibles a las 
características y respuestas de cada usuario. La literatura sobre
adaptación basada en la experiencia del jugador ha identificado 
que aún persiste una brecha entre la captura de información del 
usuario, la interpretación de su estado y la aplicación de 
cambios efectivos en tiempo real del sistema \cite{lopes2025closing}.
Esto es significativo porque la personalización no depende únicamente 
de observar métricas, sino de convertirlas en decisiones capaces de 
ajustar la experiencia de manera coherente con las necesidades del 
jugador. En este sentido, una experiencia musical interactiva representa 
un escenario particularmente valioso, ya que en ella no solo importa el 
cumplimiento de objetivos o tareas, sino también la conservación de un 
equilibrio entre reto, inmersión y apreciación de la música y del 
entorno \cite{pichlmair2007sound}. Precisamente por esa necesidad de 
balance, este tipo de propuesta ofrece un contexto adecuado para 
examinar el potencial de la adaptación dinámica como medio para 
enriquecer la experiencia del usuario.

Desde el punto de vista técnico, la adaptación dinámica puede 
construirse a partir de variables observables del comportamiento del 
jugador. En particular, se ha reportado que sistemas de ajuste dinámico 
de dificultad pueden modificar parámetros del juego en tiempo real con 
base en el desempeño del usuario, por ejemplo, mediante cambios en 
velocidad, complejidad o estructura de las tareas \cite{huber2021dda}.
De igual forma, revisiones recientes sobre estimación de 
\textit{engagement} en videojuegos indican que las métricas utilizadas 
con mayor frecuencia para modelar el estado del jugador incluyen señales 
de desempeño, patrones de interacción, telemetría del juego y, en ciertos 
casos, información fisiológica, lo que confirma que la captura de 
métricas en \textit{runtime} constituyen una base válida para construir 
mecanismos de adaptación \cite{rashed2025engagement}.

La implementación de adaptación dinámica también presenta dificultades 
que justifican la necesidad de diseñar cuidadosamente la relación entre 
métricas y modificaciones del \textit{gameplay} 
\cite{xue2017dda,zohaib2018dda}. Un sistema excesivamente sensible puede 
responder a variaciones aisladas del desempeño y producir cambios 
frecuentes o contradictorios en la dificultad. De igual forma, 
modificaciones demasiado evidentes pueden afectar la percepción de 
control del jugador o alterar la comprensión de las reglas de la 
experiencia \cite{lopes2025closing,zook2012temporal}. Por ello, la
adaptación no consiste únicamente en detectar un error y disminuir 
inmediatamente la dificultad, sino en interpretar tendencias 
suficientemente estables antes de modificar las condiciones del juego.

Además de controlar cuándo y con qué frecuencia se realizan los ajustes, 
debe considerarse que un buen desempeño no representa necesariamente una 
experiencia satisfactoria, así como la presencia de errores no implica 
por sí misma una experiencia negativa \cite{yannakakis2009adaptation,rashed2025engagement}.
Esta consideración adquiere mayor importancia en VR, donde constructos 
como presencia, experiencia de usuario y \textit{cybersickness} son 
dimensiones centrales para valorar la calidad de una aplicación 
inmersiva \cite{wang2025presence,bareisyte2024questionnaires}. 
Incrementar factores como la velocidad, cantidad de objetos o la 
complejidad puede elevar simultáneamente las exigencias de la 
actividad \cite{sepich2022workload}. El análisis de estos constructos se 
beneficia de la combinación entre métricas objetivas y herramientas 
subjetivas de evaluación. También se ha observado que el 
\textit{feedback} visual y háptico puede incrementar la sensación de 
presencia, mientras que la carga de tarea puede influir en mareos y otros 
malestares \cite{sepich2022workload,gibbs2022feedback}. Por ello, resulta 
relevante diseñar un sistema que utilice varias métricas de interacción 
y reglas de adaptación controladas, procurando conservar un equilibrio 
entre desafío, claridad, \textit{confort} e involucramiento en lugar de 
optimizar únicamente la cantidad de aciertos del usuario.

Otro valor que tiene este proyecto radica en la necesidad de estudiar 
las experiencias inmersivas desde una perspectiva más integral, ya que 
diversos componentes que influyen en la calidad de la interacción, como 
la retroalimentación multisensorial, el comportamiento del usuario y la 
adaptación del sistema, suelen analizarse de manera aislada \cite{chen2024interfaces,spittle2023interaction,yeo2024interfaces,lopes2025closing,huber2021dda}.
Esta separación limita la comprensión de cómo dichos elementos se 
relacionan entre sí dentro de una experiencia real de uso. En 
consecuencia, resulta pertinente desarrollar propuestas que permitan 
observar de forma conjunta cómo la interacción, las métricas del 
usuario y la adaptación dinámica pueden incidir en la personalización y 
en la capacidad de respuesta del sistema frente al contexto de uso \cite{lopes2025closing,rashed2025engagement}.
Visto así, el proyecto adquiere valor tanto para la investigación como 
para la industria, al aportar un caso aplicado que puede contribuir a 
la comprensión de cómo diseñar experiencias gamificadas más sensibles 
al usuario y mejor fundamentadas en datos de interacción.